令 $\Aff_k^\pf$ 表示完全 $k$-代数构成的范畴。我们通过 $\Aff_k^\pf$ 上的预层来定义 $G$ 在 $k$ 上的仿射 Grassmann 簇 $\Gr_G$，
\begin{equation}
\label{E:affGrass def}
\Gr_G(R)= \left\{(\mE,\beta)\left | \begin{array}{l}\mE \mbox{ 是 } D_R\mbox{ 上的 } G\mbox{-torsor，且 }  \\
 \beta: \mE\dashrightarrow \mE^0 \mbox{ 是一个修正}\end{array}\right.\right\},
\end{equation}
其中 $\mE^0$ 表示平凡 $G$-torsor。若群 $G$ 明确，为简便起见把 $\Gr_G$ 写作 $\Gr$。
\begin{thm}\label{affGrass}
仿射 Grassmann 簇 $\Gr$ 可由子函子的归纳极限表示：$\Gr=\underrightarrow\lim X_i$，其中 $X_i\to X_{i+1}$ 为闭浸入，$X_i$ 是射影簇的完全化。
\end{thm}
当 $\cha F=\cha k$ 时，这是经典结果，归功于 Beauville–Laszlo、Faltings 等人（参看 \cite{BL, Fa}）。在混合特征情形（即 $\cha F\neq \cha k$），这作为猜想出现于 \cite[附录 B]{Zh1}，并由 Bhatt–Scholze 证明（\cite{BS}）。此前 \cite{Zh1} 证明了较弱的结论：$X_i$ 是紧合代数空间的完全化；这已足以在此设定下表述几何 Satake。

\begin{rmk}
完全 $k$-概形范畴是完全 $k$-代数上预层范畴的全子范畴，见 \cite[引理 A.10]{Zh1}。故上述定理有意义。
\end{rmk}
\begin{rmk}读者或许会问，我们为何只把仿射 Grassmann 簇 $\Gr$ 看作 $\Aff_k^\pf$ 上的预层，而不像通常那样看作 $k$-代数全范畴 $\Aff_k$ 上的预层。在等特征情形这确实可行，也正是 \cite{BL, Fa} 所考虑的设定。但
在混合特征情形，由于对非完全的特征 $p$ 环取 Witt 向量环存在某些病态行为，我们不知道 $\Gr$ 作为 $\Aff_k$ 上预层的正确刻画。另一方面，过渡到完全化不改变概形的拓扑，故若只关心与 $\Gr$ 有关的拓扑问题（即 $\Gr$ 上的反常层），我们不会丢失任何信息。
\end{rmk}
\begin{rmk}
\label{R:scholzediamond}
若 $F=k((t))$，则存在定义在圆盘 $D$ 上的 Beilinson–Drinfeld Grassmann 簇（事实上，Beilinson–Drinfeld Grassmann 簇定义在 $D$ 关于 $k$ 的所有笛卡尔幂之上），使得上述定义的 $\Gr$ 恰是闭点 $s\in D$ 处的纤维。在混合特征情形，Scholze 在发明称为 diamond 的新几何结构之后，近期构造了类似的几何对象（参看 \cite{SW}）。借助这些对象，像 $\Spec \bQ_p\times \Spec \bQ_p$ 这样的纤维积才真正有意义，也才可以在这些纤维积上定义类似的 Beilinson–Drinfeld Grassmann 簇。
\end{rmk}

仿射 Grassmann 簇还有一个有用的解释，即把它看作齐性空间。
设 $H$ 为 $\mO$ 上有限型的仿射群概形。以 $L^+H$（相应地 $LH$）记 $H$ 的\textbf{射流群}（相应地\textbf{环路群}）。作为预层，
\begin{equation*}\label{E:jet and loop}
L^+H(R)=H(W_\mO(R)),\quad LH(R)=H(W_\mO(R)[1/\varpi]).
\end{equation*}
它可由一个仿射群概形（相应地归纳概形）表示。对 $r\geq 0$，令 $L^rH$ 为第 $r$ 个射流群，即
\begin{equation*}\label{E:rjet}
L^rH(R)=H(W_{\mO,r}(R)).
\end{equation*} 
则 $L^rH$ 可由某个代数 $k$-群的完全化表示（通常的 Greenberg 实现），且 $L^+H=\underleftarrow\lim L^rH$。

现在可在 $\Gr$ 上定义空间 $\Gr^{(\infty)}$，它由 $(\mE,\beta)\in \Gr(R)$ 连同平凡化 $\epsilon:\mE\simeq \mE^0$ 构成。可以证明，它是 $\Gr$ 上的 $L^+G$-torsor，且有典范同构 $\Gr^{(\infty)}\cong LG$（见 \cite[引理 1.3]{Zh1} 或 \cite[命题 1.3.6]{Zh2}）。因此可把 $\Gr$ 实现为 fpqc 商
\begin{equation}
\label{E:gr as quotient}
\Gr\cong LG/L^+G.
\end{equation}
由此 $L^+G$ 通过左乘自然作用在 $\Gr$ 上。就模空间问题 \eqref{E:affGrass def} 而言，这一作用即 $\beta$ 与自同构 $L^+G=\Aut(\mE^0)$ 的自然复合。
\medskip

 
下面定义 $\Gr$ 的一些闭子簇，例 \ref{minu Sch} 与例 \ref{quasiminuSch} 所研究的对象作为特例包含在内。
首先设 $\mE_1$ 与 $\mE_2$ 是 $D_{\bar k}=\Spec \mO_L$ 上两个 $G$-torsor，$\beta: \mE_1\dashrightarrow \mE_2$ 为修正。
我们按如下方式给 $\beta$ 配一个元素
$G(\mO_L)\backslash G(L)/G(\mO_L)$：选取同构 $\epsilon_1:\mE_1\simeq \mE^0$ 与 $\epsilon_2:\mE_2\simeq \mE^0$，便得到平凡 $G$-torsor 的自同构 $\epsilon_2\beta\epsilon_1^{-1}\in \Aut(\mE^0|_{D_{\bar k}^*})$，从而得到 $G(L)$ 中一个元素。$\epsilon_1$ 与 $\epsilon_2$ 的另选会把此元素左乘、右乘 $G(\mO_L)$ 中的元素。因此，
$\beta$ 给出 $G(\mO_L)\backslash G(L)/G(\mO_L)$ 中一个良定义的元素。经由双射 \eqref{cartan2}，给 $\beta$ 配一个元素
\[\inv(\beta)\in \xcoch(T)/W.\]
我们称 $\inv(\beta)$ 为 $\beta$ 的相对位置。
\begin{ex}\label{invbetaGLn}
设 $G=\GL_n$。按通常方式把 $D$ 上的 $\GL_n$-torsor 与秩 $n$ 的有限投射 $\mO$-模等同。则 $\inv(\beta)$ 可如下描述。
如例 \ref{coweightGLn} 那样等同 $\xcoch(T)\simeq \bZ^n$。
给定秩 $n$ 的两个有限自由 $\mO$-模 $\La_1$ 与 $\La_2$，以及 $\beta: \La_1[1/\varpi]\simeq \La_2[1/\varpi]$，总存在 $\La_1$ 的基 $(e_1,\ldots,e_n)$ 与 $\La_2$ 的基 $(f_1,\ldots,f_n)$，使 $\beta$ 由
$$\beta(e_i)=p^{m_i}f_i$$ 
给出，且 $m_1\geq m_2\geq\cdots\geq m_n$。此外，序列 $(m_1,\ldots,m_n)$ 与基的选取无关。于是定义 $\inv(\beta)=(m_1,\ldots,m_n)$。
\end{ex}
